\documentclass{situationist-preprint}

\title{Shoemaker's Frozen Worlds}
\msid{TSM-17-b}
\msfamily{TSM-17}
\msdiscipline{Time, Space and Modality}
\mstype{Consequence}
\mskeywords{time without change; Shoemaker; relationism; temporal metric; duration}

\begin{document}
\maketitle

\begin{abstract}
TSM-17-a argues, following Shoemaker, that time without change is possible and could be known to occur. We derive three consequences for the worlds Shoemaker describes. First, the relationist who holds that time is constituted by change, or by relations among changes, must deny that such worlds are possible, and is refuted. Second, the duration of an interval is not a measure of the change within it, since a total freeze has a determinate length of one year. Third, the frozen intervals are stretches of time within the B-series, not gaps between times or times outside it.
\end{abstract}

\section{Three consequences}

In the world of \citet{Shoemaker1969}, three regions freeze for a year every third, fourth and fifth year respectively, and the simplest law fitting the observations implies a total freeze every sixtieth year. By Theorem~3.3 of TSM-17-a, such a total freeze is possible and could be known to occur.

\begin{corollary}\label{c:rel}
Reductive relationism about time, on which there is nothing to time beyond the order of changes, is false.
\end{corollary}

\begin{proof}
On reductive relationism, two changes with no change between them are adjacent in time, so no interval elapses during a total freeze. In Shoemaker's world an interval elapses: the inhabitants have reason to believe, and may know, that a year passes between the last change before the joint freeze and the first change after it. Since the world is possible (Proposition~3.1 of TSM-17-a), reductive relationism is not necessarily true, and as an account of what time is it is false.
\end{proof}

\begin{corollary}\label{c:metric}
The duration of an interval is not determined by the amount of change within it.
\end{corollary}

\begin{proof}
A total freeze contains no change and has a duration of one year, the same as a local freeze in one region and as a year without any freeze. Intervals that agree in duration differ in change, and intervals that differ in change agree in duration.
\end{proof}

Corollary~\ref{c:metric} rejects Aristotle's account of time as the number of change in respect of before and after \citep[IV.11]{Aristotle}, and it bears on the metric: the inhabitants know the length of the freeze because the laws fix the periodicity of local freezes in units of time, not in units of change \citep{NewtonSmith1980}.

\begin{corollary}\label{c:within}
A total freeze is a stretch of time between its boundaries, ordered by earlier and later with the times around it.
\end{corollary}

\begin{proof}
The freeze begins after the last change before it and ends before the first change after it, and it has a duration (Corollary~\ref{c:metric}). It is therefore an interval in the order of times, as TSM-01 takes time to be: an order by earlier and later, with no need of change to be populated.
\end{proof}

Corollary~\ref{c:within} is what TSM-14-a requires: a frozen interval is not a time outside time, and its possibility does not make the question of what preceded the beginning of time well formed.

\begin{objection}
If a total freeze has no causal consequences, its duration is idle, and the inhabitants could posit a freeze of any length.
\end{objection}

\begin{reply}
The duration is fixed by the same law that implies the freeze, and the law is confirmed by the observed local freezes \citep{Warmbrod2004}. A rival hypothesis assigning another length is an unmotivated equivalent of the kind TSM-17-a shows does not defeat justification.
\end{reply}

\section{Conclusion}

Shoemaker's frozen worlds refute reductive relationism, separate the measure of time from the measure of change, and locate empty intervals within the order of times. Time without change is possible, and in such worlds it has a determinate length.

\begin{thebibliography}{99}\small
\bibitem[Aristotle()]{Aristotle} Aristotle. \emph{Physics}. Trans. R.~P. Hardie and R.~K. Gaye, in J. Barnes (ed.), \emph{The Complete Works of Aristotle}, Princeton University Press, 1984.
\bibitem[Newton-Smith(1980)]{NewtonSmith1980} Newton-Smith, W.~H. (1980). \emph{The Structure of Time}. Routledge \& Kegan Paul.
\bibitem[Shoemaker(1969)]{Shoemaker1969} Shoemaker, S. (1969). Time without change. \emph{The Journal of Philosophy}, 66(12), 363--381.
\bibitem[Warmbrod(2004)]{Warmbrod2004} Warmbrod, K. (2004). Temporal vacua. \emph{The Philosophical Quarterly}, 54(215), 266--286.
\end{thebibliography}

\end{document}
